\documentclass{beamer}
\beamertemplatenavigationsymbolsempty

% movie15 is an older package for video inclusion in Beamer.
% This file exercises the most common options exposed by \includemovie
% to verify that each one is translated correctly into PowerPoint.
\usepackage{movie15}

\title{Beamer with Video (movie15)}
\author{Author Name}
\date{\today}

\begin{document}

\begin{frame}
    \titlepage
\end{frame}

\begin{frame}{Default --- Click to Play}
    No special options: the video sits behind its poster image and PowerPoint
    starts playback when the slide is clicked.
    \begin{center}
        \includemovie[
            poster=image.png,
            text={Click to play}
        ]{0.7\linewidth}{0.4\linewidth}{video.mp4}
    \end{center}
\end{frame}

\begin{frame}{Autoplay}
    The \texttt{autoplay} option starts playback automatically when the slide
    becomes visible.
    \begin{center}
        \includemovie[
            poster=image.png,
            autoplay
        ]{0.7\linewidth}{0.4\linewidth}{video.mp4}
    \end{center}
\end{frame}

\begin{frame}{Autoplay + Repeat (loop)}
    The \texttt{repeat} option causes the video to loop indefinitely.
    Combined with \texttt{autoplay} it plays continuously while the slide is shown.
    \begin{center}
        \includemovie[
            poster=image.png,
            autoplay,
            repeat
        ]{0.7\linewidth}{0.4\linewidth}{video.mp4}
    \end{center}
\end{frame}

\begin{frame}{Autoplay + Reduced Volume (30\%)}
    The \texttt{volume} option (range 0.0--1.0) sets initial playback volume.
    \begin{center}
        \includemovie[
            poster=image.png,
            autoplay,
            volume=0.3
        ]{0.7\linewidth}{0.4\linewidth}{video.mp4}
    \end{center}
\end{frame}

\end{document}